\section*{अध्याय \ref{ch:b1:intermolecular-forces-solvents} --- अंतराआण्विक बल और विलायक}

\begin{solution}{exo:b1:intermolecular-forces-solvents:1}
आर्गन: केवल लंदन। हाइड्रोजन क्लोराइड: कीसोम, डिबाई और लंदन, जिनमें लंदन
प्रधान है (\cref{exo:b1:intermolecular-forces-solvents:8}); क्लोरीन
प्रबल \omterm{def:b1:intermolecular-forces-solvents:hbond}{हाइड्रोजन आबंधों} के लिए बहुत बड़ा है। पानी: \omterm{def:b1:intermolecular-forces-solvents:hbond}{हाइड्रोजन आबंध} प्रधान हैं, साथ में
तीनों वान्डरवाल्स पद। मेथेन: केवल लंदन। डाइआयोडीन: केवल लंदन,
जो प्रबल है क्योंकि अणु बड़ा और अत्यधिक ध्रुवणीय है।
\end{solution}

\begin{solution}{exo:b1:intermolecular-forces-solvents:2}
अपने ही अणुओं के बीच: \ce{CH3OH} (\ce{O-H}), \ce{CH3NH2}
(\ce{N-H}) और \ce{HF}। \ce{CH3OCH3} और \ce{CH3F} में O
या F पर \omterm{def:b1:lewis-resonance-vsepr:lewis}{एकाकी युग्म} हैं पर उन पर कोई हाइड्रोजन नहीं: वे केवल \omterm{def:b1:intermolecular-forces-solvents:hbond}{हाइड्रोजन आबंध} ग्रहण कर सकते हैं,
उदाहरण के लिए पानी से। \ce{CH4} न देता है न ग्रहण करता है।
\end{solution}

\begin{solution}{exo:b1:intermolecular-forces-solvents:3}
पानी और एथेनॉल: ध्रुवीय, प्रोटिक। प्रोपेनोन और ऐसीटोनाइट्राइल: ध्रुवीय,
अप्रोटिक। डाइक्लोरोमेथेन: अल्प ध्रुवीय, अप्रोटिक। हेक्सेन: अध्रुवीय, अप्रोटिक।
\end{solution}

\begin{solution}{exo:b1:intermolecular-forces-solvents:4}
डाइआयोडीन अध्रुवीय है: हेक्सेन में वह \omterm{def:b1:intermolecular-forces-solvents:vdw}{लंदन अन्योन्यक्रियाओं} के बदले \omterm{def:b1:intermolecular-forces-solvents:vdw}{लंदन
अन्योन्यक्रियाएँ} पाता है, जबकि पानी में उसे \omterm{def:b1:intermolecular-forces-solvents:hbond}{हाइड्रोजन आबंध} तोड़ने पड़ते और उनकी
जगह कुछ न मिलता। सोडियम क्लोराइड आयनिक है: केवल उच्च
परावैद्युतांक वाला विलायक उसके आयनों को अलग कर सकता है और केवल \omterm{def:b1:intermolecular-forces-solvents:solvent-classes}{ध्रुवीय विलायक} उनका \omterm{def:b1:intermolecular-forces-solvents:dissolution}{विलायकन};
हेक्सेन ($\varepsilon_r = 1.89$) दोनों में से कुछ नहीं करता।
\end{solution}

\begin{solution}{exo:b1:intermolecular-forces-solvents:5}
मेथॉक्सीमेथेन में कोई \ce{O-H} नहीं: वह \omterm{def:b1:intermolecular-forces-solvents:hbond}{हाइड्रोजन आबंध} दे नहीं सकता और सबसे
नीचे उबलता है। मेथेनॉल और एथेनॉल दोनों हाइड्रोजन-आबंधित हैं; एथेनॉल में एक
\ce{CH2} अधिक है, इसलिए उसकी \omterm{def:b1:intermolecular-forces-solvents:vdw}{लंदन अन्योन्यक्रियाएँ} प्रबलतर हैं। एथेनोइक अम्ल
हाइड्रोजन-आबंधित द्वितय बनाता है और सबसे भारी है: वह सबसे ऊपर उबलता है।
\end{solution}

\begin{solution}{exo:b1:intermolecular-forces-solvents:6}
निर्वात में $F = e^2/(4\pi\varepsilon_0 r^2) = (1.602 \times
10^{-19})^2/(4\pi \times 8.854 \times 10^{-12} \times (5.00 \times
10^{-10})^2) = \qty{9.23e-10}{N}$। पानी में $F/78.4 = \qty{1.18e-11}{N}$;
हेक्सेन में $F/1.89 = \qty{4.88e-10}{N}$। पानी आकर्षण को हेक्सेन से 41
गुना अधिक दुर्बल करता है: वह आयनों को अलग रख सकता है, हेक्सेन नहीं।
\end{solution}

\begin{solution}{exo:b1:intermolecular-forces-solvents:7}
प्रोपेनोन को पानी के साथ मिलाने से पानी--पानी के कुछ \omterm{def:b1:intermolecular-forces-solvents:hbond}{हाइड्रोजन आबंध} टूटते हैं, पर
पानी और \ce{C=O} के ऑक्सीजन के बीच नए बनते हैं, साथ में
द्विध्रुव--द्विध्रुव अन्योन्यक्रियाएँ भी: संतुलन अनुकूल है। हेक्सेन पानी को
उन \omterm{def:b1:intermolecular-forces-solvents:hbond}{हाइड्रोजन आबंधों} के बदले, जो टूटते, केवल दुर्बल \omterm{def:b1:intermolecular-forces-solvents:vdw}{लंदन अन्योन्यक्रियाएँ} देता है:
दोनों द्रव अलग रहते हैं।
\end{solution}

\begin{solution}{exo:b1:intermolecular-forces-solvents:8}
\ce{HBr} में अधिक इलेक्ट्रॉन हैं और उसका मेघ बड़ा, अधिक ध्रुवणीय है: \omterm{def:b1:intermolecular-forces-solvents:vdw}{लंदन
अन्योन्यक्रिया} में \ce{HCl} पर उसकी बढ़त, \omterm{def:b1:intermolecular-forces-solvents:vdw}{कीसोम अन्योन्यक्रिया} में उसकी कमी से
अधिक है। लंदन प्रधान है।
\end{solution}

\begin{solution}{exo:b1:intermolecular-forces-solvents:9}
\chemfig{CH_3-C(=[:60]O)-[:-60]O-[:0]H} अपने साथी के सामने, दो
$\ce{O-H}\cdots\ce{O=C}$ आबंध आठ सदस्यों वाला वलय बंद करते हैं। अध्रुवीय और
अप्रोटिक बेंज़ीन में इन आबंधों से कोई प्रतिस्पर्धा नहीं करता: अम्ल
द्वितय बना रहता है और लगभग \qty{120}{g/mol} के कण की तरह व्यवहार करता है। पानी में
हर अम्ल अणु इसके बजाय पानी के साथ \omterm{def:b1:intermolecular-forces-solvents:hbond}{हाइड्रोजन आबंध} बनाता है: द्वितय टूट जाते हैं।
\end{solution}

\begin{solution}{exo:b1:intermolecular-forces-solvents:10}
सोडियम क्लोराइड आयनों से बना है: उच्च परावैद्युतांक वाला पानी उन्हें
अलग करता है (वियोजन) और उनका जलयोजन करता है (ऑक्सीजन \ce{Na+} की ओर, हाइड्रोजन
\ce{Cl-} की ओर)। हाइड्रोजन क्लोराइड एक \omterm{def:b1:lewis-resonance-vsepr:dipole}{ध्रुवीय अणु} है: ध्रुवीय पानी
\ce{H-Cl} आबंध को आयन-युग्म \ce{H3O+ Cl-} में बदलता है (आयनन),
युग्म को अलग करता है (वियोजन) और आयनों का जलयोजन करता है। हेक्सेन अध्रुवीय है
और कम परावैद्युतांक वाला: वह न आयनन करता है न वियोजन, इसलिए धारा ढोने के लिए
कोई आयन नहीं होते।
\end{solution}

\begin{solution}{exo:b1:intermolecular-forces-solvents:11}
प्रति \ce{CH2} $(174.1 - 36.1)/5 = \qty{27.6}{\celsius}$। अंतिम
वृद्धि (नोनेन से डेकेन) \qty{23.4}{\celsius} है, इसलिए अनडेकेन लगभग
$174 + 22 \approx \qty{196}{\celsius}$ पर उबलना चाहिए। हर \ce{CH2} उतना ही
लंदन योगदान जोड़ता है, पर वह अणु के कुल योगदान का घटता हुआ
अंश होता है, जबकि अन्योन्यक्रियाओं पर विजय पाने के लिए आवश्यक ताप
पहले से ही ऊँचा है।
% ledger: bp:C9H20
\end{solution}

\begin{solution}{exo:b1:intermolecular-forces-solvents:12}
\omterm{def:b1:intermolecular-forces-solvents:amphiphile}{जलरागी} भाग: सल्फ़ेट शीर्ष \ce{-OSO3^-}, अपने \ce{Na+} प्रति-आयन के साथ;
\omterm{def:b1:intermolecular-forces-solvents:amphiphile}{जलविरागी} भाग: बारह-कार्बन श्रृंखला। सतह पर शीर्ष
पानी में और पूँछें हवा में; मिसेल में पूँछें भीतर और शीर्ष बाहर।
पूँछें धब्बे की चिकनाई में घुल जाती हैं, शीर्ष उसे पानी के संपर्क में
रखते हैं, और चिकनाई मिसेलों में उठकर अलग हो जाती है।
\end{solution}

\begin{solution}{pb:b1:intermolecular-forces-solvents:1}
\textbf{1.} \omterm{def:b1:intermolecular-forces-solvents:vdw}{लंदन अन्योन्यक्रिया}, अध्रुवीय परमाणुओं के बीच एकमात्र अन्योन्यक्रिया।
\textbf{2.} हीलियम से ज़ेनॉन तक इलेक्ट्रॉनों की संख्या और आकार बढ़ते हैं,
अतः \omterm{def:b1:periodicity:polarisability}{ध्रुवणीयता} और लंदन आकर्षण भी।
\textbf{3.} हीलियम के दो इलेक्ट्रॉनों को उसका नाभिक कसकर थामता है: यह
सबसे कम ध्रुवणीय परमाणु है, सबसे दुर्बल लंदन आकर्षण के साथ।
\textbf{4.} $68.8 - 36.1 = \qty{32.7}{\celsius}$: एक \ce{CH2}
इलेक्ट्रॉन और संपर्क-सतह जोड़ता है, अतः लंदन आकर्षण भी।
\textbf{5.} लंदन, जो \omterm{def:b1:periodicity:polarisability}{ध्रुवणीयता} के साथ बढ़ता है: $\ce{CH4} < \ce{SiH4} <
\ce{GeH4}$।
\textbf{6.} कार्बन इतना विद्युत-ऋणात्मक नहीं कि \ce{C-H} को
प्रबलता से ध्रुवित करे, और मेथेन के पास \omterm{def:b1:intermolecular-forces-solvents:hbond}{हाइड्रोजन आबंध} ग्रहण करने के लिए कोई \omterm{def:b1:lewis-resonance-vsepr:lewis}{एकाकी युग्म} नहीं।
\textbf{7.} पानी: ध्रुवीय, प्रोटिक, सबसे उच्च परावैद्युतांक वाला; यह
दोनों आयनों को अलग करता है और उनका \omterm{def:b1:intermolecular-forces-solvents:dissolution}{विलायकन} करता है।
\textbf{8.} $78.4/1.89 = 41$: पानी में आकर्षण 41 गुना दुर्बल है।
\textbf{9.} ऐसीटोनाइट्राइल (या प्रोपेनोन): लवण को घोलने लायक ध्रुवीय,
और अप्रोटिक, ताकि ऋणायन मुक्त रहे।
\textbf{10.} हेक्सेन: पानी के साथ अमिश्रणीय, डाइआयोडीन की तरह अध्रुवीय। आयोडीन
हेक्सेन में, ऊपरी परत में, चला जाता है (हेक्सेन पानी से कम घना है)।
\textbf{11.} एथेनॉल पानी के साथ \omterm{def:b1:intermolecular-forces-solvents:miscibility}{मिश्रणीय} है: दूसरी परत बनती ही नहीं।
\textbf{12.} इसका ऑक्सीजन पानी से \omterm{def:b1:intermolecular-forces-solvents:hbond}{हाइड्रोजन आबंध} ग्रहण करता है, पर यह स्वयं
कोई आबंध नहीं दे सकता और इसके दोनों एथिल समूह \omterm{def:b1:intermolecular-forces-solvents:amphiphile}{जलविरागी} हैं: यह केवल
थोड़ा घुलता है।
\textbf{13.} प्रति \omterm{def:b1:periodicity:table}{आवर्त} $-66.4 - (-85.0) = \qty{18.6}{\celsius}$;
बहिर्वेशित \ce{HF}: $-85.0 - 18.6 = \qty{-103.6}{\celsius}$।
\textbf{14.} \ce{HF} \qty{19.6}{\celsius} पर उबलता है, बहिर्वेशन से
\qty{123}{\celsius} ऊपर।
\textbf{15.} $-62.5 - (-87.8) = \qty{25.3}{\celsius}$; बहिर्वेशित
\ce{NH3}: \qty{-113.1}{\celsius}; वास्तविक $-33.4$, \qty{80}{\celsius}
अधिक।
\textbf{16.} $-88.1 - (-112.0) = \qty{23.9}{\celsius}$; बहिर्वेशित
\ce{CH4}: \qty{-135.9}{\celsius}; वास्तविक $-161.5$, नीचे। मेथेन कोई
\omterm{def:b1:intermolecular-forces-solvents:hbond}{हाइड्रोजन आबंध} नहीं बनाता; वह बस छोटा और कम ध्रुवणीय है।
\textbf{17.} पानी (लगभग \qty{179}{\celsius}) $>$ \ce{HF} (123) $>$
\ce{NH3} (80)।
\textbf{18.} \ce{NH3} तीन दे सकता है पर केवल एक ग्रहण कर सकता है (एक \omterm{def:b1:lewis-resonance-vsepr:lewis}{एकाकी युग्म});
\ce{HF} तीन ग्रहण कर सकता है पर केवल एक दे सकता है। पानी दो देता है और दो ग्रहण
करता है: हर अणु चार \omterm{def:b1:intermolecular-forces-solvents:hbond}{हाइड्रोजन आबंधों} में भाग ले सकता है, एक पूरा
त्रिविमीय जाल।
\textbf{19.} \ce{H2S} और \ce{H2Se} कोई उल्लेखनीय \omterm{def:b1:intermolecular-forces-solvents:hbond}{हाइड्रोजन आबंध} नहीं बनाते:
उनके क्वथनांक उस लंदन प्रवृत्ति का पालन करते हैं जिसका पालन पानी
उनके बिना करता। स्वयं पानी ही वह विसंगति है जिसे मापा जा रहा है।
\textbf{20.} प्रति \omterm{def:b1:periodicity:table}{आवर्त} $-41.3 - (-60.3) = \qty{19.0}{\celsius}$।
\textbf{21.} $T(p) = -60.3 + 19.0\,(p - 3)$, \unit{\celsius} में।
\textbf{22.} $T(2) = -60.3 - 19.0 = \qty{-79.3}{\celsius}$।
\textbf{23.} $100.0 - (-79.3) = \qty{179}{\celsius}$।
\textbf{24.} अपने \omterm{def:b1:intermolecular-forces-solvents:hbond}{हाइड्रोजन आबंधों} के बिना पानी लगभग
\textbf{\qty{-79}{\celsius}} पर उबलता, और पृथ्वी पर गैस होता।
\end{solution}
